\RequirePackage{iftex}
\ifPDFTeX
\documentclass[a4paper,12pt]{article}
\usepackage[margin=24mm]{geometry}
\usepackage[T1]{fontenc}
\usepackage[utf8]{inputenc}
\usepackage{graphicx}
\usepackage{CJKutf8}
\else
\documentclass[a4paper,12pt]{jsarticle}
\usepackage[dvipdfmx,margin=24mm]{geometry}
\usepackage[dvipdfmx]{graphicx}
\fi
\usepackage{amsmath,amssymb}
\usepackage{enumitem}

\setlength{\parindent}{0pt}
\setlength{\parskip}{0.35em}
\setlist[enumerate]{leftmargin=3em,itemsep=0.25em,topsep=0.35em}

\newcommand{\tenpoints}{\hfill （10点）}
\newcommand{\fivepoints}{\hfill （5点）}
\newcommand{\subtenpoints}{\hfill （10点）}
\newcommand{\blankbox}[2]{\fbox{\rule{0pt}{2.4ex}\makebox[#1]{#2}}}
\newcommand{\problem}[2]{\par\vspace{0.9em}\noindent{\large 問題 #1}\quad（#2点）\quad\ignorespaces}

\begin{document}

\ifPDFTeX
\begin{CJK}{UTF8}{min}
\fi

\begin{center}
{\Large 数学1A レポート1}　満点 100点
\end{center}

計算問題では途中式も書くこと。証明問題では，どの定理を使ったか書くこと。

\problem{1}{$5\times 2=10$}
次の文を，$\mathbb{R},\mathbb{N},\forall,\exists$ を用いて記号で書き直せ。

\begin{enumerate}[label=(\alph*)]
\item 任意の実数 $x$ に対して，$x$ より大きい自然数 $n$ が存在する。 
\item ある実数 $a$ が存在して， $a^2 = 2 $ となる。 
\end{enumerate}

\problem{2}{10}
$f(x)=x^2$ とする。このとき
\(
\lim_{x\to 1} f(x)=1
\) 
であることを，次の四角を埋めることで示せ:  $\varepsilon>0$ を任意にとる。このとき, 
\[
\delta=\blankbox{10em}{A}
\]
とおく。$0<|x-1|<\delta$ のとき，
\[
|f(x)-1|
=|x^2-1|
=|x-1||x+1|
<\blankbox{8em}{B}
<\varepsilon.
\]

\problem{3}{10}
赤道上の各地点の気温を角度 $\theta$ の関数 $T(\theta)$ で表す。$T$ は連続であり，$T(\theta+2\pi)=T(\theta)$ を満たすと仮定する。 このとき，赤道上には互いに反対側にある2地点で，気温が等しいものが存在することを中間値の定理を用いて示せ。

\problem{4}{$5\times 2=10$}
次の関数の微分を求めよ。

(1) $\displaystyle  (x^2+1)\sin x$ 
(2) $\displaystyle \frac{e^x}{1+x^2}$ 

\problem{5}{10}
$0<a<b$ とする。平均値の定理を用いて，次を示せ。

\[
\frac{\sin b-\sin a}{b-a}\leq 1
\]

\problem{6}{$5\times 2=10$}
ロピタルの定理を用いて，次の極限を求めよ。

\begin{enumerate}[label=(\alph*)]
\item $\displaystyle \lim_{x\to 0}\frac{\sin x}{x}$ 
\item $\displaystyle \lim_{x\to 0}\frac{e^x-1-x}{x^2}$ 
\end{enumerate}

\problem{7}{$5\times 2=10$}
次の関数の高階微分を求めよ。

\begin{enumerate}[label=(\alph*)]
\item $f(x)=x^2e^x$ の $f^{(3)}(x)$ 
\item $g(x)=\sin x$ の $g^{(2026)}(x)$ 
\end{enumerate}

\problem{8}{$10\times 2=20$}
次の関数の マクローリン展開を求めよ。

(1) $\sin x$  (2) $\log(1+x)$

\problem{9}{10}
テイラー展開の歴史について調べ，最低５行以上でまとめよ。

\ifPDFTeX
\end{CJK}
\fi

\end{document}
